\loai{Loại 1. Nhận nhiệt, nhận công}
\begin{vidu}%[1P6-3.2B][Loai1] 
Người ta thực hiện một công 100 J lên một khối khí và truyền cho khối khí một nhiệt lượng 40 J. Độ biến thiên nội năng của khí là
\choice
{60 J và nội năng giảm}
{\True 140 J và nội năng tăng}
{60 J và nội năng tăng}
{140 J và nội năng giảm}
\loigiai{ 
\begin{itemize}
\item Khí nhận nhiệt: $Q=40\ J$.
\item Thực hiện công lên khối khí suy ra khí nhận công: $A=100\ J$.
\item Độ biến thiên nội năng của khí: $\Delta U=A+Q=140\ J>0\Rightarrow$ nội năng của khí tăng 140 J.
\end{itemize}}
\end{vidu}	
\loai{Loại 2. Nhận nhiệt, sinh công}
\begin{vidu}%[1P6-3.2B][Loai1] 
Người ta truyền cho khí trong xi-lanh nhiệt lượng 100 J. Khí nở ra thực hiện công 70 J đẩy pit-tông lên. Độ biến thiên nội năng của khí là
\choice
{$-30$ J}
{170 J}
{\True 30 J}
{$-170$ J}
\loigiai{
\begin{itemize}
\item Khí nhận nhiệt: $Q=100\ J$.
\item Khí dãn nở sinh công: $A=-70\ J$.
\item Độ biến thiên nội năng của khí: $\Delta U=A+Q=30\ J$.
\end{itemize}}
\end{vidu}
\loai{Loại 3. Nhận công, tỏa nhiệt}
\begin{vidu}%[1P6-3.2B][Loai1] 
Thực hiện công 100 J để nén khí trong xi-lanh và khí truyền ra môi trường một nhiệt lượng 20 J thì
\choice
{\True nội năng của khí tăng 80 J}
{nội năng của khí tăng 120 J}
{nội năng của khí giảm 80 J}
{nội năng của khí giảm 120 J}
\loigiai{
\Binhluan{
\begin{itemize}
\item Khí truyền nhiệt: $Q=-20\ J$.
\item Nén khí, thể tích giảm nên khí nhận công: $A=100\ J$.
\item Độ biến thiên nội năng của khí: $\Delta U=A+Q=80\ J>0$. 
\item Nội năng của khí tăng 80 J.
\end{itemize}
}
{"Ta" thực hiện công nên "khối khí" nhận công: $A>0$}
}
\end{vidu}
\loai{Loại 4. Tính công}
\begin{vidu}%[1P6K2-2][Loai1]
Nội năng của khối khí tăng 10 J khi truyền cho khối khí một nhiệt lượng 30 J. Khi đó khối khí đã
\choice
{\True thực hiện công là 20 J}
{nhận công là 40 J}
{nhận công là 20 J}
{sinh công là 40 J}
\loigiai{
\Binhluan{
\begin{itemize}
\item Ta có: $\Delta U = 10\ J,\ Q = 30\ J \Rightarrow A = -20\ J$. 
\item Do đó khối khí thực hiện công là 20 J.
\end{itemize}
}
{Dấu $-$ chứng tỏ khối khí đã thực hiện công}
}
\end{vidu}

\loai{Loại 5. Tính nhiệt lượng}
\begin{vidu}%[1P6-2.2K][Loai1]
Người ta thực hiện một công 170 J lên khối khí trong xi-lanh. Biết rằng nội năng của khí tăng thêm 170 J. Khi đó khối khí
\choice
{tỏa nhiệt 340 J}
{nhận nhiệt 340 J}
{nhận nhiệt 170 J}
{\True không trao đổi nhiệt với môi trường}
\loigiai{
Ta có: $A = 170\ J,\ \Delta U = 170\ J \Rightarrow Q = \Delta U-A=0 \Rightarrow$ Khối khí không trao đổi nhiệt với môi trường.
\begin{note}
Người thực hiện công lên khối khí $\Rightarrow$ khối khí nhận công $\Rightarrow A>0$.
\end{note}
}
\haidong{6}\end{vidu}
\loai{Loại 6. Tính công do khí thực hiện rồi tính độ biến thiên nội năng}
\begin{vidu}
Người ta cung cấp một nhiệt lượng $Q=10\ J$ cho khối khí trong một xi-lanh đặt nằm ngang. Khối khí dãn nở, đẩy pit-tông chuyển động thẳng đều một đoạn $0,1\ m$ và pit-tông chịu tổng lực cản có độ lớn $F_c=20\ N$.
\begin{vdc}
Công mà khí thực hiện là
\choice
{10\ J}
{\True 2\ J}
{8\ J}
{0,5\ J}
%\loigiai{
%\begin{itemize}
%\item[A] Sai. Đây là nhiệt lượng $Q$ chứ không phải công.
%\item[B] Đúng. $A = F.d = 20.0,1 = 2\ J$
%\item[C] Sai. Đây là độ biến thiên nội năng chứ không phải công.
%\item[D] Sai. Tính sai công.
%\end{itemize}
%}

\end{vdc}
\dotlineEX{4}
\begin{vdc}
Nội năng của khí đã thay đổi như thế nào?
\choice
{Giảm $2\ J$}
{Không đổi}
{\True Tăng $8\ J$}
{Tăng $12\ J$}
%\loigiai{
%\begin{itemize}
%\item[A] Sai. Khí không mất nội năng, mà nhận nhiệt lượng và sinh công.
%\item[B] Sai. Có nhiệt lượng truyền vào và công được thực hiện nên nội năng có thay đổi.
%\item[C] Đúng. $\Delta U = Q + A = 10 + (-2) = 8\ J$
%\item[D] Sai. Tính nhầm, cộng sai giá trị.
%\end{itemize}
%}

\end{vdc}
\dotlineEX{4}
\end{vidu}



\begin{vidu}
Khi truyền nhiệt lượng $Q$ cho khối khí trong một xi-lanh hình trụ thì khí dãn nở đẩy pit-tông làm thể tích của khối khí tăng thêm $7\ \ell$. Biết áp suất của khối khí là $3.10^5\ Pa$ và không đổi trong quá trình khí dãn nở. Trong quá trình này, nội năng của khối khí tăng $1100\ J$.
\begin{vdc}
Độ lớn công mà khí thực hiện là
\choice
{$2,1.10^6\ J$}
{\True $2100\ J$}
{$-2100\ J$}
{$4,3.10^7\ J$}
%\loigiai{
%\begin{itemize}
%\item[A] Sai: Sai đơn vị. Không đổi đơn vị từ lít ra $m^3$ nên kết quả quá lớn.
%\item[B] Đúng: $A = p.\Delta V = 3.10^5.\dfrac{7}{1000} = 2100\ J$
%\item[C] Sai: Độ lớn công là số dương; hướng công âm không ảnh hưởng đến độ lớn.
%\item[D] Sai: Phép chia này vô lý và không đúng về mặt đại số và vật lý.
%\end{itemize}
%}

\end{vdc}
\dotlineEX{4}
\begin{vdc}
Nhiệt lượng $Q$ đã truyền cho khí có giá trị là
\choice
{$-1000\ J$}
{\True $3200\ J$}
{$1000\ J$}
{$2300\ J$}
%\loigiai{
%\begin{itemize}
%\item[A] Sai: Đổi dấu công sai, công thực hiện là dương nên không lấy dấu âm.
%\item[B] Đúng: $\mathcal{Q} = \Delta U + A = 1100 + 2100 = 3200\ J$
%\item[C] Sai: Công thức sai, công không trừ nội năng.
%\item[D] Sai: Mặc dù đúng phép cộng nhưng giải thích sai bản chất: $Q = \Delta U + A$, không phải đảo ngược.
%\end{itemize}
%}

\end{vdc}
\dotlineEX{4}
\end{vidu}



\begin{vidu}
Một chất khí đựng trong bình hình trụ được lắp một pit-tông có thể chuyển động không ma sát trong bình. Khi hấp thụ nhiệt lượng $400 \ J$ thì chất khí trong bình giãn nở dưới áp suất bên ngoài không đổi là $1,00$ atm từ thể tích $5,00$ lít đến $10,0$ lít. Cho biết $1\ \ell .atm$ tương đương với $101,3 \ J$. Độ biến thiên nội năng của khí trong bình là bao nhiêu jun (làm tròn kết quả đến chữ số hàng đơn vị)?
\shortans[oly]{-107}
\loigiai{
\begin{itemize}
\item Công thực hiện bởi khí khi giãn nở:
\[ A = P \Delta V = 1,00 \text{atm} . (10,0 \text{lít} - 5,00 \text{lít}) . 101,3 \text{J/atm} \]
\[ A = 1,00 . 5,00 . 101,3 = 506,5\ J \]

\item Theo định luật I của nhiệt động lực học:
\[ Q = \Delta U + A \]
\[ 400 = \Delta U + 506,5 \]
\[ \Delta U = 400 - 506,5 = -106,5\ J \]				
\end{itemize}	
}
\end{vidu}
\begin{vidu}
Một lượng khí chứa trong một xi-lanh đặt thẳng đứng có pit-tông di chuyển được. Ở trạng thái cân bằng, chất khí chiếm thể tích $V\left(m^3\right)$ và tác dụng lên pit-tông một áp suất $4.10^{5} \ N/m^2$. Khối khí nhận một nhiệt lượng $1000 \ J$ và giãn nở đẩy pit-tông đi lên thẳng đều làm thể tích khí tăng thêm $0,003\ m^3$. Coi rằng áp suất chất khí không đổi. 
\choiceTF[t]
{\True Lượng khí bên trong xi-lanh nhận nhiệt và sinh công làm biến đổi nội năng}
{Theo quy ước, khối khí nhận nhiệt và sinh công nên $A>0;\ Q>0$}
{\True Công mà khối khí thực hiện có độ lớn bằng $1200 \ J$}
{\True Độ biến thiên nội năng của khối khí $\Delta U=-200 \ J$}
\loigiai{
\begin{enumerate}[label=\alph*)]
\item ĐÚNG
\item SAI. Khối khí nhận nhiệt và sinh công nên $Q>0,\ A<0$. 
\item ĐÚNG. $|A|=p.\Delta V=4.10^5.0,003=1200\ J$.
\item ĐÚNG. $\Delta U=A+Q=-1200+1000=-200\ J$.
\end{enumerate}
}
\haidong{7}\end{vidu}
\begin{vidu} \nguon{3}
Một lượng khí trong một xi-lanh hình trụ bị nung nóng, khí nở ra đẩy pit-tông lên làm thể tích tăng thêm $0,02\ m^3$ và nội năng tăng thêm $1280\ J$. Biết áp suất của khối khí là $2 . 10^{5}\ Pa$ và không đổi trong quá trình dãn nở. Nhiệt lượng đã truyền cho khí bằng bao nhiêu jun?
\shortans[oly]{5280 } 
\loigiai{
\begin{itemize}
\item Độ lớn công của khối khí thực hiện: $A^{\prime}=F . \ell=p S . \ell=p \Delta V=2 . 10^{5} . 0,02=4000\ J$
\item Công do khí sinh ra nên ta phải viết: $A=-A^{\prime}=-4000\ J$
\item Theo nguyên lí I NĐLH: $\Delta U=A+Q \Rightarrow 1280=-4000+Q \Rightarrow Q=5280\ J$
\end{itemize}
}
\end{vidu}